\documentclass[conference]{IEEEtran}
\usepackage{cite}
\usepackage{graphicx}
\usepackage{url}
\usepackage{amsmath}
\usepackage{booktabs}

\begin{document}

\title{Challenges of Generative AI in STEM Education}

\author{
\IEEEauthorblockN{Petronela Vieru}
\IEEEauthorblockA{
Politehnica University of Timisoara\\
Timisoara, Romania\\
petronela.vieru@student.upt.ro}
}

\maketitle

\begin{abstract}
This paper investigates the challenges associated with the adoption of Generative Artificial Intelligence (GenAI) in STEM education. This document is a template intended for further development.
\end{abstract}

\begin{IEEEkeywords}
Generative AI, STEM Education, Artificial Intelligence, Educational Technology
\end{IEEEkeywords}

\section{Introduction}
% Add background, motivation and objectives.

\section{Literature Review}
\subsection{Generative AI in Education}
% Add literature review.

\subsection{Generative AI in STEM Education}
% Add literature review.

\section{Research Methodology}
% Describe methodology.

\section{Challenges of Generative AI in STEM Education}
\subsection{Academic Integrity}
\subsection{Reliability and Accuracy}
\subsection{Ethical and Privacy Concerns}
\subsection{Dependence on AI Tools}

\section{Discussion}
% Add discussion.

\section{Future Research Directions}
% Add future work.

\section{Conclusion}
% Add conclusions.

\bibliographystyle{IEEEtran}
\bibliography{references}

\end{document}